\section{考点总结}
\begin{longtable}{P{2.6cm}P{4.3cm}P{4.1cm}P{3cm}}
\toprule 模型 & 适用条件 & 必须检查 & 对应题目\\\midrule\endhead
前缀和 & 静态多次区域和 & 0号哨兵、容斥符号、整数范围 & 邻域均值\\
前后缀汇总 & 单点变化影响一整段 & 变化前后的分界是否包含当前位置 & 梦境巡查\\
差分 & 修改全部累积后再恢复 & r+1边界、恢复的成本 & 画图\\
二分答案 & 判定沿答案单调 & 证明单调、哨兵、区间缩小、乘法 & 机器人宿管\\
双指针 & 有序性或窗口单调性 & 安全移动理由、重复值、总移动次数 & 买菜、稀疏向量\\
贪心 & 有交换论证或最优前缀性质 & 反例、每次选择能否固定 & 分蛋糕\\
质因数分解 & 试除数量可承受 & 完整指数、残留质数、n变小 & 因子化简\\
搜索 & 状态完整、规模允许 & 去重、撤销、剪枝不能漏解 & 机器人复健\\
动态规划 & 子问题可复用 & 状态、转移、初值、顺序、答案位置 & 机器人饲养指南\\
STL & 操作与数据性质吻合 & 严格比较、空容器、失效、复杂度 & 数字排序、相似度计算\\
\bottomrule\end{longtable}
\textbf{复盘要求：}给每题写三句话：为什么直接枚举会慢（或其实够用）；利用了哪个性质；若输入去掉排序、正数、静态或恰好装满等条件，原算法哪一步会失效。能回答这些问题，比记住代码模板更有用。
